\section{Riemannian Diffusion Models}
\label{sec:riemannian_diffusion_models}
Riemannian diffusion models (RDMs) generalise Euclidean diffusion by defining
the generative process on a smooth Riemannian manifold $(\mathcal{M},g)$, where
$g$ is a metric that determines local distances, angles, curvature, and
geodesic structure. This setting is appropriate when the intrinsic structure of
the data lies on a curved or constrained manifold—such as anatomical surfaces,
pose manifolds, molecular conformations, or latent spaces with non-linear
geometry—where Euclidean diffusion would distort semantic neighbourhoods or
produce artefacts. RDMs replace Euclidean Brownian motion with manifold
Brownian motion, governed by the \emph{Laplace–Beltrami operator}
$\Delta_{\mathcal{M}}$, and adapt the reverse process accordingly. This yields
a generative model whose trajectories respect the intrinsic geometry of the
data.

\paragraph{Mathematical Foundations.}

Let $(\mathcal{M}, g)$ be a $d$-dimensional Riemannian manifold with metric
tensor $g$, volume measure $\mathrm{d}\mathrm{vol}_g$, and geodesic distance
$d_{\mathcal{M}}(\cdot,\cdot)$. Riemannian diffusion replaces the Euclidean SDE
\[
\mathrm{d}x_t
    = f_t(x_t)\,\mathrm{d}t + g_t \,\mathrm{d}w_t
\]
with a manifold-valued SDE whose noise lives in the tangent space
$T_x\mathcal{M}$:
\[
\mathrm{d}x_t
    = V_t(x_t)\,\mathrm{d}t + \sqrt{2}\, \mathrm{d}W_t^{\mathcal{M}},
\]
where $W_t^{\mathcal{M}}$ is Brownian motion on $\mathcal{M}$ and whose
generator is the Laplace--Beltrami operator
\[
\Delta_{\mathcal{M}} = \mathrm{div}_g \circ \nabla^{\mathcal{M}}.
\]

The probability density $p_t$ evolves according to the \emph{Riemannian heat
equation}
\[
\partial_t p_t = \Delta_{\mathcal{M}} p_t.
\]
The score is defined via the Riemannian gradient
\[
\nabla^{\mathcal{M}} \log p_t
    = g^{-1} \nabla_{\text{coord}} \log p_t,
\]
and the reverse-time SDE (Anderson, Song et al.\ adapted to $\mathcal{M}$) is
\[
\mathrm{d}x_t
    = \left( V_t(x_t)
        - 2\, \nabla^{\mathcal{M}} \log p_t(x_t)
      \right) \mathrm{d}t
      + \sqrt{2}\,\mathrm{d}\bar{W}_t^{\mathcal{M}}.
\]
Thus the score network in RDMs estimates a \emph{Riemannian} gradient instead of
a Euclidean one, ensuring that sampling trajectories never leave the manifold.

\paragraph{Riemannian ELBO and Heat Kernel Construction.}

The Riemannian Diffusion Models framework introduces a continuous-time ELBO
using a Riemannian Feynman–Kac representation:
\[
\log p(x_0)
    \ge \mathbb{E}_{x_{0:T}}\!\left[
        \log p_T(x_T)
        + \int_0^T
            \left\langle
                \nabla^{\mathcal{M}} \log p_t(x_t),
                \mathrm{d}W_t^{\mathcal{M}}
            \right\rangle_g
        \right],
\]
where $p_T$ is a Riemannian Gaussian defined via the heat kernel
\[
K_t(x,y)
    \approx (4\pi t)^{-d/2}
        \exp\!\left(
            -\frac{d_{\mathcal{M}}(x,y)^2}{4t}
        \right).
\]
This formulation provides exact training objectives for diffusion on curved
spaces and clarifies how curvature influences the score and drift fields.

\paragraph{Practical Realisation and Scalability.}

The ``Scaling Riemannian Diffusion Models'' work addresses the computational
difficulty of computing exact geometric operators. It introduces:
\begin{itemize}
    \item \emph{Neural metric learning}: learning $g(x)$ directly from data.
    \item \emph{Local coordinate charts}: approximating $\mathcal{M}$ via small
          Euclidean patches with transition consistency.
    \item \emph{Approximate geodesic solvers}: computing
          $d_{\mathcal{M}}(x,y)$ and exponential maps efficiently.
    \item \emph{Efficient Riemannian score matching}: replacing costly
          Laplace–Beltrami operations with tractable approximations.
\end{itemize}
These innovations make RDMs trainable with complexity comparable to standard
diffusion, enabling application to high-dimensional data including medical
images.

\paragraph{Why Riemannian Geometry Matters for Compositionality.}

Diffusion is inherently geometric: sampling follows trajectories governed by
the score, whose structure reflects the curvature and topology of the underlying
data manifold. For compositional generative modelling, the following insights
are crucial:
\begin{itemize}
    \item \textbf{H1 (Latent Alignment):} Two models can only be composed
          meaningfully when their manifolds have aligned tangent spaces and
          compatible coordinate charts. Misalignment causes inconsistent geodesic
          directions and broken semantic fusion.
    \item \textbf{H2 (Local Disentanglement):} Semantic factors correspond to
          directions in $T_x \mathcal{M}$. If these are entangled or curved
          incompatibly across models, semantic composition fails even if density
          composition succeeds.
    \item \textbf{H4 (Curvature Compatibility):} If the manifolds underlying two
          models have incompatible curvature (e.g., one flat, one highly curved),
          the combined log-density geometry becomes inconsistent, producing
          anatomically implausible hybrids.
\end{itemize}

\paragraph{Implications for This Research.}

Medical images, including pathological variations, often lie on structured
submanifolds shaped by anatomical and physiological constraints. TB, Silicosis,
and Silico–TB differ not only in texture and intensity but also in geometric
structure: cavity morphology, nodule distribution, and pleural interactions
define distinct submanifolds or manifold branches. Riemannian diffusion provides
precise mathematical tools—geodesic drift alignment, curvature estimation,
metric learning, and Laplace–Beltrami analysis—to determine when two pretrained
models inhabit compatible geometric spaces and when composition will succeed or
fail.

These tools directly support our Phase~3 experiments, where we analyse geometric
compatibility and use Riemannian principles to explain and predict the behaviour
of Euclidean compositional methods (score-based, PoE, SuperDiff). The
Riemannian framework thus offers the theoretical depth required to interpret the
structural prerequisites underlying semantically meaningful generative
composition.
